% !tp-example Basic usage
\documentclass{article}
\begin{document}
\wcTpPythonViewer[src=/docs/components/_shared/intro-projects/python.json]{}
\end{document}

% !tp-example Javascript API
\documentclass{article}
\begin{document}
\wcScript[type=module]|
import "/components/python/python.js";
import "/components/python-viewer/python-viewer.js";

const example = document.createDocumentFragment();

const element1 = document.createElement("tp-python-viewer");
element1.setAttribute("src", "/docs/components/_shared/intro-projects/python.json");
example.append(element1);

document.body.append(example);
|
\end{document}

% !tp-example src
\documentclass{article}
\begin{document}
\wcTpPythonViewer[src=/docs/components/python-viewer/examples/example.json]{}
\end{document}

% !tp-example script
\documentclass{article}
\begin{document}
\wcTpPythonViewer{
  \wcScript[type=tp/python]|
        message = "Hello Python viewer"
        print(message)
      |
}
\end{document}

% !tp-example matplotlib
\documentclass{article}
\begin{document}
\wcTpPythonViewer[src=/docs/components/python-viewer/examples/matplotlib.json]{}
\end{document}

% !tp-example error
\documentclass{article}
\begin{document}
\wcTpPythonViewer{
  \wcScript[type=tp/python]|
        raise RuntimeError("Intentional Python error")
      |
}
\end{document}

% !tp-example Single source file
\documentclass{article}
\begin{document}
\wcTpPythonViewer[src=/docs/components/_shared/single-source/example.py]{}
\end{document}

% !tp-example Compare presentation modes
\documentclass{article}
\begin{document}
\wcSection{
  \subsection*{mode="default" — Standard controls}
  Open Code to edit the source, then click Run or press \wcKbd*{Shift}+\wcKbd*{Enter} in the editor.

  Run to print a console message followed by an intentional error. In simple mode the console opens automatically. Comment out the message statements and run again to check that the error alone opens it; remove the error statement to try a successful run.

  \wcTpPythonViewer[mode=default]{
    \wcScript[type=tp/python,filename=main.py]|
from js import document
document.body.innerHTML = "<p>The program ran.</p>"
print("Console message before the intentional error.")
raise RuntimeError("Intentional error: the simple-mode console opens automatically.")
  |
  }
}

\wcSection{
  \subsection*{mode="lite" — Compact controls}
  Open Code to edit the source, then click Run or press \wcKbd*{Shift}+\wcKbd*{Enter} in the editor.

  Run to print a console message followed by an intentional error. In simple mode the console opens automatically. Comment out the message statements and run again to check that the error alone opens it; remove the error statement to try a successful run.

  \wcTpPythonViewer[mode=lite]{
    \wcScript[type=tp/python,filename=main.py]|
from js import document
document.body.innerHTML = "<p>The program ran.</p>"
print("Console message before the intentional error.")
raise RuntimeError("Intentional error: the simple-mode console opens automatically.")
  |
  }
}

\wcSection{
  \subsection*{mode="simple" — Editor and rendered output only}
  Edit the code below, keep the cursor in the editor and press \wcKbd*{Shift}+\wcKbd*{Enter} (Return) to run it. This viewer has no Run button.

  Run to print a console message followed by an intentional error. In simple mode the console opens automatically. Comment out the message statements and run again to check that the error alone opens it; remove the error statement to try a successful run.

  \wcTpPythonViewer[mode=simple]{
    \wcScript[type=tp/python,filename=main.py]|
from js import document
document.body.innerHTML = "<p>The program ran.</p>"
print("Console message before the intentional error.")
raise RuntimeError("Intentional error: the simple-mode console opens automatically.")
  |
  }
}

Programming viewers already have compact controls, so default and lite look alike. The first run may load the language engine.
\end{document}
